\section{Universal Contracts: Adapter, Benchmark, and CertOS Governance}
\label{sec:prof-runtime-contracts}

The professor-facing ORIUS draft should present three contracts as a single
stack: the adapter contract, the benchmark contract, and the CertOS
governance contract.  Together they explain how degraded observation is
converted into a typed runtime safety layer rather than a battery-only
control recipe.

\subsection{Adapter contract}
The adapter is the only place where domain-specific telemetry and
constraints are allowed to enter the universal runtime.  The current
monograph and supporting design documents make this boundary explicit:
\texttt{src/orius/certos/runtime.py} provides a domain-neutral governance
surface, \texttt{docs/UNIVERSAL\_BENCHMARK\_SPEC.md} defines the common
step schema, and \texttt{docs/UNIVERSAL\_GOVERNANCE\_SPEC.md} requires a
domain governance policy instead of battery-shaped runtime state.

In practical terms, each domain adapter must expose a true-state violation
predicate, an observed-state satisfiability predicate, an optional margin
computation, a fallback action, and certificate-compatible action repair.
That is the minimum needed to make ORIUS portable across battery energy
storage, autonomous vehicles, industrial control, healthcare monitoring,
navigation, and aerospace without changing the meaning of safety at the
kernel level.

\subsection{Benchmark contract}
ORIUS-Bench is the manuscript-facing replay contract. The universal step
record keeps the same fields across domains:
\texttt{true\_constraint\_violated},
\texttt{observed\_constraint\_satisfied},
\texttt{true\_margin},
\texttt{observed\_margin},
\texttt{intervened},
\texttt{fallback\_used},
\texttt{certificate\_valid},
\texttt{latency\_us}, and
\texttt{domain\_metrics}. The metrics family is equally fixed: TSVR, OASG,
certificate validity accuracy, intervention rate, audit completeness,
latency percentiles, and coverage diagnostics when conformal uncertainty is
active.

This contract matters because it prevents the professor-facing draft from
collapsing the whole manuscript into a single battery result table.  It
also keeps the evidence boundary honest: the same schema can be reused
across all six rows, but the rows do not carry the same proof depth.

\subsection{CertOS governance contract}
CertOS is the runtime governance layer.  Its role is deliberately narrow:
do not release an action unless the runtime can either validate a current
certificate or replace the action with an explicit fallback.  The policy
surface is domain-neutral and consists of five hooks:
\texttt{is\_actuation(action)},
\texttt{fallback\_action(constraints, state)},
\texttt{horizon\_update(certificate, reliability, drift, constraints, validity\_horizon)},
\texttt{required\_certificate\_fields()}, and
\texttt{integrity\_projection(action)}.

The invariants are correspondingly small and auditable: no actuation
without a valid or fallback certificate path, an unbroken certificate hash
chain, and fallback only when validity collapses. That is the right level
of abstraction for a professor-facing draft that wants to argue for a
fundamental safety layer.

\begin{algorithm}[t]
\caption{CertOS Certificate Release and Lifecycle Update}
\label{alg:prof-certos}
\begin{algorithmic}[1]
\Require repaired action \(a_t^{\mathrm{safe}}\), certificate candidate \(c_t\), reliability \(w_t\), constraints \(\kappa_t\)
\State project \(a_t^{\mathrm{safe}}\) through \texttt{integrity\_projection}
\If{\texttt{required\_certificate\_fields()} missing}
  \State mark certificate invalid and request fallback
\EndIf
\State validate lifecycle state and audit hash continuity
\If{certificate valid}
  \State compute horizon update using \texttt{horizon\_update}
  \State append validated certificate to audit ledger
  \State release projected action
\Else
  \State replace with \texttt{fallback\_action} and append degraded certificate
  \State release fallback only if lifecycle invariants still hold
\EndIf
\end{algorithmic}
\end{algorithm}

The professor-facing version should end this section with the claim
boundary stated plainly: ORIUS is a universal runtime safety contract for
Physical AI under degraded observation, but the evidence remains tiered by
domain and should not be presented as equal-depth proof across all rows.
